\SourceLeaf{267}
\Article{1}
\indent Von der Aufgabe:
\begin{quote}
die Theile einer gegebenen Fläche auf einer andern gegebenen Fläche so abzubilden, dass die Abbildung dem abgebildeten in den kleinsten Theilen ähnlich wird
\end{quote}
\noindent habe ich im Jahre 1822 eine allgemeine Auflösung gegeben, welche Herr Conferenzrath \textsc{Schumacher} im 3 Heft der Astronomischen Abhandlungen hat abdrucken lassen. Bei der Anwendung dieser Aufgabe auf die höhere Geodäsie, für welche sie eine vorzüglich ergiebige Hülfsquelle wird, macht sich das Bedürfniss fühlbar, Abbildungen, welche unter der angegebenen Bedingung stehen, durch eine besondere Benennung auszuzeichnen, und ich werde daher dieselben \textit{conforme} Abbildungen oder Übertragungen nennen, indem ich diesem sonst vagen Beiworte eine mathematisch scharf bestimmte Bedeutung beilege.

In der angeführten Schrift ist die allgemeine Auflösung, welche eine willkürliche Function einschliesst, auch auf mehrere \textit{bestimmte} Flächen angewandt; das letzte dort behandelte Beispiel betrifft die conforme Übertragung der Oberfläche des Umdrehungsellipsoids auf die Kugelfläche, und es ist [Art. 13] zugleich eine solche Bestimmung der arbiträren Function angegeben, die zu einer sehr brauchbaren Anwendung auf die höhere Geodäsie benutzt werden kann. Diese Benutzung war a. a. O. nur kurz angedeutet, und eine ausführlichere Entwickelung vorbehalten. Ich werde jedoch anstatt \textit{dieser} speciellen Auflösung eine etwas abgeänderte und für die geodätischen Anwendungen noch viel mehr geeignete Methode zur conformen Übertragung der ellipsoidischen Fläche auf die Kugelfläche in der gegenwärtigen Abhandlung entwickeln, und damit zugleich alles zu einer solchen Benutzung erforderliche verbinden.

\Article{2}
\indent Die allgemeine Auflösung der Aufgabe, angewandt auf die ellipsoidische und sphärische Fläche, gibt folgende alle conformen Übertragungen der einen auf die andere umfassende Formel (1):
\[
T+i\log\operatorname{cotg}\tfrac{1}{2}U=f(t+i\log\{\operatorname{cotg}\tfrac{1}{2}w\mathord{.}(\tfrac{1-e\cos w}{1+e\cos w})^{\frac{1}{2}e}\})
\]

Es bezeichnen hier
